% SP1 comprimido para la v2: una figura de composición, las tres ecuaciones
% que certifican cada etapa, el enunciado de cada resultado formal (la
% demostración completa va al Anexo~\ref{app:sp1-full-detail}) y una tabla de
% resultados por etapa. La etapa E3 ya estaba en este formato en la v1
% (sections/sp1-wrench-condensed.tex) y se reutiliza sin cambios.
\subsection{SP1: formación distribuida de coaliciones}
\label{subsec:sp1}
\zlabel{budget:sp1-start}
% Intervalos y valores p recalculados con la SEMILLA como unidad de
% remuestreo por final-hardening/recompute_clustered_ci.py.
% Bootstrap pareado sobre semillas, B=10000, semilla de analisis 20260919;
% Wilcoxon de una cola en la direccion predeclarada sobre diferencias
% por semilla, y Holm dentro de la familia de la campana. Los
% contrastes deterministas por semilla no llevan intervalo ni valor p.
\newcommand{\SPTwoWorlds}{1560}
\newcommand{\SPTwoComparisonRows}{20280}
\newcommand{\SPTwoAblationRows}{9360}
\newcommand{\SPTwoScoreCoverage}{0.514}
\newcommand{\SPTwoScoreSuccess}{0.487}
\newcommand{\SPTwoScoreRuntime}{40.84}
\newcommand{\SPTwoCoverageCoverage}{0.529}
\newcommand{\SPTwoCoverageSuccess}{0.227}
\newcommand{\SPTwoCoverageRuntime}{11.31}
\newcommand{\SPTwoNeuralGap}{0.109}
\newcommand{\SPTwoNeuralSuccess}{0.263}
\newcommand{\SPTwoNeuralRuntime}{1.01}
\newcommand{\SPTwoNeuralParameters}{121}
\newcommand{\SPTwoSmithGap}{0.343}
\newcommand{\SPTwoSmithSuccess}{0.103}
\newcommand{\SPTwoSmithRuntime}{0.35}
\newcommand{\SPTwoPlainSuccess}{0.051}
\newcommand{\SPTwoMarginalSuccess}{0.135}
\newcommand{\SPTwoPlainGap}{0.372}
\newcommand{\SPTwoMarginalGap}{0.157}
\newcommand{\SPTwoPlainIncomplete}{0.946}
\newcommand{\SPTwoMarginalIncomplete}{0.848}
\newcommand{\SPTwoPlainAlignment}{0.049}
\newcommand{\SPTwoMarginalAlignment}{0.137}
\newcommand{\SPTwoGapEffect}{-0.2145}
\newcommand{\SPTwoGapCILow}{-0.2230}
\newcommand{\SPTwoGapCIHigh}{-0.2065}
\newcommand{\SPTwoGapPHolm}{4{,}55\times10^{-12}}
\newcommand{\SPTwoSuccessEffect}{0.0838}
\newcommand{\SPTwoSuccessCILow}{0.0763}
\newcommand{\SPTwoSuccessCIHigh}{0.0915}
\newcommand{\SPTwoIncompleteEffect}{-0.0982}
\newcommand{\SPTwoAlignmentEffect}{0.0877}
\newcommand{\SPTwoPDEffect}{-0.119}
\newcommand{\SPTwoPDCILow}{-0.127}
\newcommand{\SPTwoPDCIHigh}{-0.111}
\newcommand{\SPTwoPDPHolm}{1{,}000}


SP1 decide la composición y los contactos antes de autorizar el movimiento de
una carga en tres etapas encadenadas, cada una con un certificado
estrictamente más fuerte que la anterior. El razonamiento completo de cada
etapa, incluidas las demostraciones y la comparación de métodos, se conserva
en el Anexo~\ref{app:sp1-full-detail}.

\begin{figure}[!htbp]
  \centering
  \begin{tikzpicture}[x=1mm,y=1mm,font=\sffamily\fontsize{7.0}{7.6}\selectfont]
    \tikzset{stage/.style={draw=VIUGray!70,line width=0.5pt,rounded corners=2pt,
        minimum width=42mm,minimum height=17mm,align=center,fill=white},
      arr/.style={-{Latex[length=2.2mm]},line width=0.7pt,draw=VIUDark!70}}
    \node[stage,fill=VIUOrange!8] (e1) at (24,10)
      {\textbf{E1 · Cuotas}\\[1mm]{\fontsize{6.0}{6.5}\selectfont certifica: cardinalidad exacta}\\
       {\fontsize{6.0}{6.5}\selectfont\color{VIUGray}no certifica: heterogeneidad}};
    \node[stage,fill=SPBlue!8] (e2) at (78,10)
      {\textbf{E2 · Servicio heterogéneo}\\[1mm]{\fontsize{6.0}{6.5}\selectfont certifica: cobertura ponderada}\\
       {\fontsize{6.0}{6.5}\selectfont\color{VIUGray}no certifica: factibilidad mecánica}};
    \node[stage,fill=SPPurple!8] (e3) at (132,10)
      {\textbf{E3 · \emph{Wrench} planar}\\[1mm]{\fontsize{6.0}{6.5}\selectfont certifica: contacto realizable}\\
       {\fontsize{6.0}{6.5}\selectfont\color{VIUGray}no certifica: estabilidad de movimiento}};
    \draw[arr] (e1.east)--(e2.west);
    \draw[arr] (e2.east)--(e3.west);
    \node[anchor=north,text=VIUGray,font=\fontsize{6.2}{6.6}\selectfont] at (78,-3)
      {SP2 certifica el movimiento a partir del contacto que entrega E3.};
  \end{tikzpicture}
  \caption[Composición de las tres etapas de SP1.]{Composición de las tres
  etapas de SP1. Cada flecha entrega un certificado que la etapa siguiente
  consume sin ampliar su dominio; ninguna etapa transfiere su garantía a la
  ejecución física, que corresponde a SP2.}
  \label{fig:sp1-compact-pipeline}
  \viuownsource
\end{figure}

\paragraph{E1 --- regulación exacta de cuotas}
Sean $\mathcal I=\{1,\ldots,N\}$, $\mathcal K=\{1,\ldots,K\}$ y $a_i\in\{0\}\cup\mathcal K$, donde $a_i=0$ deja al robot libre y $q_k(a)=|\{i:a_i=k\}|$. Para cuotas obligatorias $n_k>0$ y costes adimensionales fijos $\kappa_{ik}\in[0,\kappa_{\max}]$,
\begin{equation}
 \begin{aligned}
 D(a)&=\sum_{k\in\mathcal K}[n_k-q_k(a)]_+,&
 O(a)&=\sum_{k\in\mathcal K}[q_k(a)-n_k]_+,\\
 C(a)&=\sum_{i:a_i\ne0}\kappa_{i,a_i},&
 \Phi_1(a)&=-\{C(a)+\lambda D(a)+\lambda O(a)\}.
 \end{aligned}
 \label{eq:pre-sp1-quota-game}
\end{equation}
Con $U_i(a)=\Phi_1(a)-\Phi_1(0,a_{-i})$ la utilidad de retirarse a la
condición libre,

\begin{contribucion}{potencial exacto y regulación finita de cuotas}
\estadoaporte{demostrado bajo activación asíncrona justa (cada robot revisa
su acción infinitas veces) y continuación maximal (el proceso no se detiene
mientras exista una mejora estricta disponible); demostración completa en el
Anexo~\ref{app:sp1-full-detail}.}
\begin{teorema}[Potencial exacto y regulación de cuotas]
\label{thm:pre-sp1-exact-quota}
Para el juego de~\eqref{eq:pre-sp1-quota-game}, si $\sum_kn_k\leq N$ y $\lambda>\kappa_{\max}$, $\Phi_1$ es un potencial exacto y los equilibrios de Nash puros son exactamente los perfiles con $q_k=n_k$ para todo $k$. Toda ejecución asíncrona justa de mejor respuesta que continúe mientras exista una mejora estricta alcanza uno de esos perfiles después de, a lo sumo, $(K+1)^N-1$ cambios aceptados. Los máximos globales de $\Phi_1$ son las asignaciones de cuota exacta con coste $C$ mínimo.
\end{teorema}
La cota cuenta cambios estratégicos aceptados, no segundos ni rondas globales. No se transfiere a Smith muestreado, al cierre QR bajo escasez, a comunicación imperfecta ni a la planta mecánica.
\end{contribucion}

\paragraph{E2 --- servicio operacional heterogéneo}
Dos coaliciones de igual cardinalidad pueden producir servicio distinto. Con $c_i^{\mathrm{pay}}$ la carga útil nominal, $b_i$ la fracción de batería y $\delta_{ik}$ la distancia inicial,
\begin{equation}
  a_{ik}=\underbrace{\max\{0,\tfrac{b_i-b_i^{\mathrm{res}}}{1-b_i^{\mathrm{res}}}\}}_{\psi_i^b}\;
  \underbrace{\exp(-\delta_{ik}/\ell_d)}_{\psi_{ik}^d},
  \qquad
  e_{ik}=\frac{c_i^{\mathrm{pay}}}{c^{\mathrm{ref}}}a_{ik},
  \qquad
  S_k(x)=\sum_ie_{ik}x_{ik}\geq d_k^{\mathrm{srv}}.
  \label{eq:sp1-compact-service}
\end{equation}
\begin{figure}[!htbp]
  \centering
  \begin{spdiagram}
    \spdrawarena{E2 \textbar\ misma cardinalidad, distinto servicio}
    \node[spload,minimum width=1.55cm,minimum height=0.88cm] (load) at (6.65, 2.75) {$L_k$};
    \node[spannotation] at (6.65, 1.78) {$d_k^{\mathrm{srv}}=10$};
    \node[sprobot] (r1) at (1.35, 4.10) {$R_1$};
    \node[sprobot] (r2) at (2.85, 3.00) {$R_2$};
    \node[spannotation,anchor=west] at (0.55, 3.45) {$e_{1k}=4$};
    \node[spannotation,anchor=west] at (0.55, 2.30) {$e_{2k}=4$};
    \draw[spcandidate] (r1) -- (load);
    \draw[spcandidate] (r2) -- (load);
    \node[spannotation,text=SPRed] at (3.10, 1.25) {$C_k^A:\ 4+4=8<10$};
    \node[sprobotgreen] (r3) at (10.45, 4.10) {$R_3$};
    \node[sprobotgreen] (r4) at (11.95, 3.00) {$R_4$};
    \node[spannotation,anchor=east] at (12.75, 3.45) {$e_{3k}=7$};
    \node[spannotation,anchor=east] at (12.75, 2.30) {$e_{4k}=5$};
    \draw[spassign] (r3) -- (load);
    \draw[spassign] (r4) -- (load);
    \node[spannotation,text=SPGreen!70!black] at (10.40, 1.25) {$C_k^B:\ 7+5=12\geq10$};
  \end{spdiagram}
  \caption[Caso didáctico E2: misma cardinalidad, distinto servicio.]{Caso didáctico E2 de la Ecuación~\eqref{eq:sp1-compact-service}: cardinalidad dos idéntica, servicio distinto por heterogeneidad de carga útil, batería y distancia; los valores no son observaciones experimentales.}
  \label{fig:sp1-compact-e2-scenario}
  \viuownsource
\end{figure}

\begin{contribucion}{puntuación de servicio y potencial marginal de E2}
\estadoaporte{demostrado para preferencias continuas y parámetros congelados; demostración completa en el Anexo~\ref{app:sp1-full-detail}.}
\begin{proposicion}[Alineación marginal del servicio operacional]
\label{prop:sp2-marginal-alignment}
En una región donde $S_k(\rho)<d_k^{\mathrm{srv}}$, con $E$, $d_k^{\mathrm{srv}}$
y los pesos fijos, la puntuación marginal corregida por la fracción de
demanda $e_{ik}/d_k^{\mathrm{srv}}$ es el gradiente exacto de un potencial
escalar. La puntuación plana, sin ese factor, no lo es en general cuando
existen $i,j,k$ con $e_{ik}\neq e_{jk}$.
\end{proposicion}
\end{contribucion}

\begin{table}[!htbp]
  \centering
  \caption{Métodos comparados en E2.}
  \label{tab:sp1-compact-e2-methods}
  \fontsize{7.6}{8.6}\selectfont
  \setlength{\tabcolsep}{2.2pt}
  \renewcommand{\arraystretch}{1.08}
  \begin{tabularx}{\textwidth}{>{\raggedright\arraybackslash}p{2.6cm} >{\raggedright\arraybackslash}p{1.9cm} Y}
    \toprule
    Método & Coste & Garantía o límite \\
    \midrule
    MILP de puntuación & exponencial en peor caso & certificado de optimalidad con brecha cero; prioriza completitud \\
    MILP de cobertura & exponencial en peor caso & techo certificado para cobertura normalizada \\
    Húngaro expandido \parencite{kuhn1955Hungarian} & $O(M^3)$ & exacto para el LAP expandido; capacidad aproximada \\
    Voraz de capacidad & $O(NK\log NK)$ & selección dependiente del orden; sin cota \\
    CBBA-capacidad \parencite{choiBrunetHow2009CBBA} & $O(NK)$ + consenso & garantía CBBA original inaplicable \\
    Replicator/BNN/Smith \parencite{sandholm2010population} & $O(NK)$ & ordenación finita; garantías de la EDO fuera de esta implementación \\
    Primal--dual & $O(NK)$ & cierre entero heurístico; sin certificado de la solución entera \\
    Imitación / red neuronal & $O(NKd)$ o $O(NKP)$ & entrenadas por imitación supervisada del oráculo MILP; evaluación empírica, sin certificado combinatorio \\
    Asignador secuencial + puntuación marginal (TFM) & $O(NK)$ & mejora la regla secuencial propia; no se combinó con los métodos anteriores \\
    \bottomrule
  \end{tabularx}
  \viusource{elaboración propia a partir de los métodos ejecutados y las fuentes citadas; detalle completo en el Anexo~\ref{app:sp1-full-detail}}
\end{table}

\begin{table}[!htbp]
  \centering
  \caption{Ablación pareada de las puntuaciones plana y marginal en E2 ($\SPTwoWorlds$ mundos).}
  \label{tab:sp1-compact-e2-ablation}
  \resizebox{0.82\textwidth}{!}{\begin{tabular}{lrrrr}
\toprule
Variante & Completitud & Brecha & Serv. incompleto & Alineación \\
\midrule
Puntuación plana & 0.051 & 0.372 & 0.946 & 0.049 \\
Puntuación marginal & 0.135 & 0.157 & 0.848 & 0.137 \\
\bottomrule
\end{tabular}
}
  \viusource{elaboración propia a partir de los datos procesados de E2}
\end{table}

La corrección marginal mejora la regla secuencial propia
($\SPTwoPlainSuccess\to\SPTwoMarginalSuccess$ de completitud), pero en
términos absolutos sigue por debajo de la red neuronal no certificada
($\SPTwoNeuralSuccess$) y del oráculo MILP ($\SPTwoScoreSuccess$). Esa
comparación con la red neuronal no es arquitectónicamente simétrica: la red se
entrena por imitación supervisada de las etiquetas del propio oráculo, de modo
que su cercanía al techo es en parte consecuencia de esa supervisión y no una
propiedad del método distribuido. No hay fuga de datos ---entrenamiento en las
semillas $1200$--$1223$, validación en $2200$--$2209$ y evaluación en
$3100$--$3139$, conjuntos disjuntos---, pero la asimetría de información se
mantiene: la regla distribuida nunca ve una etiqueta del oráculo. No se evaluó
combinada con los métodos aproximados de la
Tabla~\ref{tab:sp1-compact-e2-methods}.
El umbral $S_k\geq d_k^{\mathrm{srv}}$ puede cumplirse sin que los robots
tengan puntos de contacto capaces de equilibrar el torque. Ese contraejemplo
motiva el certificado de E3. El Anexo~\ref{app:sp1-full-detail} desarrolla las
cifras, los intervalos y los ocho métodos completos.

\subsubsection{Etapa 1.4: certificado planar de contacto}
\label{subsubsec:sp1-contacts}

% Intervalos y valores p recalculados con la SEMILLA como unidad de
% remuestreo por final-hardening/recompute_clustered_ci.py.
% Bootstrap pareado sobre semillas, B=10000, semilla de analisis 20260919;
% Wilcoxon de una cola en la direccion predeclarada sobre diferencias
% por semilla, y Holm dentro de la familia de la campana. Los
% contrastes deterministas por semilla no llevan intervalo ni valor p.
\newcommand{\SPThreeWorlds}{600}
\newcommand{\SPThreeRuns}{7200}
\newcommand{\SPThreeGuardedPrecision}{1.000}
\newcommand{\SPThreeGuardedCoverage}{1.000}
\newcommand{\SPThreeGuardedAbstention}{0.500}
\newcommand{\SPThreeGuardedFP}{0.000}
\newcommand{\SPThreeUnguardedFP}{0.333}
\newcommand{\SPThreeGuardEffect}{-0.333}
\newcommand{\SPThreeGuardedGap}{0.00018}
\newcommand{\SPThreeScalarGap}{0.09351}
\newcommand{\SPThreeVectorEffect}{-0.09333}
\newcommand{\SPThreeVectorCILow}{-0.09359}
\newcommand{\SPThreeVectorCIHigh}{-0.09307}
\newcommand{\SPThreePairGap}{0.00001}
\newcommand{\SPThreeCBBAGap}{0.09335}
\newcommand{\SPThreePairEffect}{-0.09334}
\newcommand{\SPThreeExactKKT}{0.00221}
\newcommand{\SPThreeRingKKT}{0.02029}
\newcommand{\SPThreeGraphEffect}{-0.01809}
\newcommand{\SPThreeGraphCILow}{-0.01907}
\newcommand{\SPThreeGraphCIHigh}{-0.01711}


E3 reemplaza la cobertura escalar por un certificado planar de puestos y
\emph{wrench}. La memoria retiene el reparto, su guardia y el contraste; la
monografía conserva el juego continuo, las derivaciones y las variantes. El
modelo es cuasiestático y anterior al transporte.

Para $Q_W\succ0$, $\varepsilon_\lambda>0$ y el producto de conjuntos de
contacto $\Lambda_k=\prod_{i\in C_k}\Lambda_{ik}$ cerrado, convexo y no vacío,
el reparto implementado resuelve
\begin{align}
 \boldsymbol\lambda_k^\star&\in\argmin_{\boldsymbol\lambda_k\in\Lambda_k}
 \|G_k\boldsymbol\lambda_k-\boldsymbol W_k^{\mathrm{dem}}\|_{Q_W}^{2}
 +\varepsilon_\lambda\|\boldsymbol\lambda_k\|_2^2, \label{eq:sp3-wrench}\\
 \rho_k^{W\star}&=\|G_k\boldsymbol\lambda_k^\star-
 \boldsymbol W_k^{\mathrm{dem}}\|_{Q_W}. \nonumber
\end{align}
Si el minimizador no restringido es interior, entonces
\begin{equation}
 \boldsymbol\lambda_k^\star=(G_k^{\mathsf T}Q_WG_k+\varepsilon_\lambda I)^{-1}
 G_k^{\mathsf T}Q_W\boldsymbol W_k^{\mathrm{dem}}.
 \label{eq:sp3-interior-wrench-solution}
\end{equation}

\begin{contribucion}{reparto único y certificado regularizado conservador}
\estadoaporte{demostrado para $\Lambda_k$ cerrado, convexo y no vacío.}
\begin{proposicion}[Residual puro frente a reparto regularizado]
\label{prop:pre-sp1-wrench-residual}
Si $\rho_k^{W,\min}=\min_{\boldsymbol\lambda_k\in\Lambda_k}
\|G_k\boldsymbol\lambda_k-\boldsymbol W_k^{\mathrm{dem}}\|_{Q_W}$, el
QP~\eqref{eq:sp3-wrench} posee minimizador único y
$\rho_k^{W,\min}\leq\rho_k^{W\star}$. Por tanto,
$\rho_k^{W\star}\leq\epsilon_W$ certifica una realización tolerable, pero la
conversa no vale en general.
\end{proposicion}
\end{contribucion}
El Anexo~\ref{app:sp1-wrench-proofs} demuestra la unicidad, la desigualdad y un
contraejemplo escalar. La selección continua que precede al cierre es
primal--dual; su potencial corregido y sus KKT se documentan en la monografía,
sin transferirlos al cierre entero ni a Euler muestreado.

La ejecución digital proyecta preferencias y actualiza precios de ocupación,
\begin{equation}
 \rho_i^{m+1}=P_{\Delta_i}(\rho_i^m+\Delta t\,F_i),\qquad
 \pi_a^{m+1}=[\pi_a^m+\Delta t_\pi h_a(\rho^m)]_+,
 \label{eq:sp3-digital-update}
\end{equation}
donde cada paso de la Ecuación~\eqref{eq:sp3-digital-update} proyecta $\rho_i$
sobre el símplex y después asigna como máximo un robot por puesto. La guardia comprueba
\eqref{eq:sp3-wrench}; el oráculo enumera instancias pequeñas, Hungarian es solo
referencia escalar y CBBA, voraz y dinámicas poblacionales son comparadores
aproximados \parencite{kuhn1955Hungarian,choiBrunetHow2009CBBA,sandholm2010population}.

\begin{table}[!htbp]
  \centering
  \caption{Contraste E3 sobre \SPThreeWorlds{} mundos pareados.}
  \label{tab:sp3-results}
  \viutablefont
  \begin{tabularx}{0.96\textwidth}{>{\raggedright\arraybackslash}p{3.65cm} C C Y}
    \toprule
    Tratamiento & Falsos positivos & Residual KKT & Lectura \\
    \midrule
    PD exacto sin guardia & \SPThreeUnguardedFP{} & \SPThreeExactKKT{} & estacionariedad no certifica cierre \\
    PD exacto con guardia & \SPThreeGuardedFP{} & \SPThreeExactKKT{} & elimina aceptaciones inviables en el banco evaluado \\
    PD anillo con guardia & \SPThreeGuardedFP{} & \SPThreeRingKKT{} & cuatro rondas dejan error de agregado \\
    \bottomrule
  \end{tabularx}
  \viuownsource
\end{table}

Las \SPThreeRuns{} ejecuciones sobre \SPThreeWorlds{} celdas pareadas, con
100~semillas predeclaradas $461000$--$461099$, redujeron los falsos positivos
de \SPThreeUnguardedFP{} a \SPThreeGuardedFP{} en las geometrías evaluadas
(efecto pareado $\SPThreeGuardEffect$). Ese promedio no es una estimación con
incertidumbre muestral: el desenlace es idéntico en las 100~semillas de cada
geometría. En \emph{off-center CoM} y \emph{slot saturation} la variante sin
guardia aceptó una asignación inviable en todas las semillas y la guardada en
ninguna; en las otras cuatro geometrías ambas coinciden en todas. El $0{,}333$
es, por tanto, la fracción de geometrías del banco donde la guardia actúa
---dos de seis---, no una tasa estimada, y publicar un intervalo o un valor $p$
para este contraste describiría la variación entre geometrías elegidas, no
variabilidad aleatoria. La consecuencia es de alcance: el resultado es
determinista dentro del banco y no se extrapola fuera de esas seis geometrías.
Tampoco acredita contacto, estabilidad ni una arquitectura completamente
vecinal; SP2 vuelve a certificar el movimiento.

\zlabel{budget:sp1-end}
